\textbf{Pour cette première partie, aucune justification n'est demandée et une seule réponse est possible par question. Pour chaque question, reportez son numéro sur votre copie et indiquez votre réponse.}

Une réponse fausse ou l'absence de réponse n'enlève aucun point.

\subsection*{Question 1}
	Un article de \np[g]{5} perd $10\,\%$ de sa masse. Sa nouvelle masse est :

	\medskip
	\begin{tblr}{colspec = {*4{X}}}
		\textbf{A.}\quad \np[g]{4,0} &
		\textbf{B.}\quad \np[g]{4,1} &
		\textbf{C.}\quad \np[g]{4,5} &
		\textbf{D.}\quad \np[g]{4,9} \\
	\end{tblr}

\subsection*{Question 2}
	Pour augmenter un prix de $15\,\%$, je dois le multiplier par :


	\begin{tblr}{colspec = {*4{X}}}
		\textbf{A.}\quad 0,15 &
		\textbf{B.}\quad 1,15 &
		\textbf{C.}\quad 0,85 &
		\textbf{D.}\quad 1,015\\
	\end{tblr}

\subsection*{Question 3}
	On considère \quad $A=\dfrac{2}{5}$. \quad Alors :

	\medskip
	\begin{tblr}{colspec = {*4{X}}}
		\textbf{A.}\quad $A = 2,5\,\%$ &
		\textbf{B.}\quad $A = 4\,\%$ &
		\textbf{C.}\quad $A = 40\,\%$ &
		\textbf{D.}\quad $A = 25\,\%$ \\
	\end{tblr}

\subsection*{Question 4}

	On considère \quad $B=\left(5^{3}\right)^{4} \times 5^{10}$. \quad Alors :

	\medskip
	\begin{tblr}{colspec = {*4{X}}}
		\textbf{A.}\quad $B=5^{17}$ &
		\textbf{B.}\quad $B=5^{120}$ &
		\textbf{C.}\quad $B=5^{22}$ &
		\textbf{D.}\quad $B=5^{70}$ \\
	\end{tblr}

\subsection*{Question 5}
	Le triple de $3^{12}$ est :

	\medskip
	\begin{tblr}{colspec = {*4{X}}}
		\textbf{A.}\quad $9^{12}$ &
		\textbf{B.}\quad $3^{36}$ &
		\textbf{C.}\quad $3^{13}$ &
		\textbf{D.}\quad $9^{36}$ \\
	\end{tblr}

\subsection*{Question 6}
	Une expression développée de $(2 x+5)^{2}$ est :

	\medskip
	\begin{tblr}{colspec = {*4{X}}}
		\textbf{A.}\quad $4 x^{2}+25$ &
		\textbf{B.}\quad $4 x^{2}+10 x+25$ &
		\textbf{C.}\quad $2 x^{2}+20 x+25$ &
		\textbf{D.}\quad $4 x^{2}+20 x+25$ \\
	\end{tblr}

\subsection*{Question 7}
	L'ensemble des solutions de l'équation $x^{2}=5$ est :

	\medskip
	\begin{tblr}{colspec = {*4{X}}}
		\textbf{A.}\quad $\big\{-2,5 ; 2,5\big\}$ &
		\textbf{B.}\quad $\big\{\sqrt{5}\big\}$ &
		\textbf{C.}\quad $\big\{2,5\big\}$ &
		\textbf{D.}\quad $\big\{-\sqrt{5} ; \sqrt{5}\big\}$
	\end{tblr}

\subsection*{Question 8}
	En convertissant $\np[mm^3]{3400}$ en $\mathrm{cm}^{3}$, on obtient :

	\medskip
	\begin{tblr}{colspec = {*4{X}}}
		\textbf{A.}\quad $\np[cm^3]{3,4}$ &
		\textbf{B.}\quad $\np[cm^3]{34}$ &
		\textbf{C.}\quad $\np[cm^3]{340}$ &
		\textbf{D.}\quad $\np[cm^3]{0,34}$
	\end{tblr}

\subsection*{Question 9}
	En sachant que le point A de coordonnées $(-3 ; y_{\mathrm{A}})$ appartient à la courbe d'équation $y=x^{2}$, on a :

	\medskip
	\begin{tblr}{colspec = {*4{X}}}
		\textbf{A.}\quad $y_{\mathrm{A}}=-9$ &
		\textbf{B.}\quad $y_{\mathrm{A}}=-6$ &
		\textbf{C.}\quad $y_{\mathrm{A}}=6$ &
		\textbf{D.}\quad $y_{\mathrm{A}}=9$
	\end{tblr}

\subsection*{Question 10}
	On lance un dé truqué à 4 faces. La probabilité d'obtenir chacune des faces est donnée par le tableau ci-dessous.

	\medskip

	\begin{tblr}{colspec={cX[c]X[c]X[c]X[c]},hlines, vlines}
		   Face     & Face numéro 1 & Face numéro 2 & Face numéro 3 & Face numéro 4 \\
		Probabilité & 0,1           & 0,25          & 0,3           & $x$
	\end{tblr}
	\smallskip

	On a alors :

	\medskip
	\begin{tblr}{colspec = {*4{X}}}
		\textbf{A.}\quad $x=0,35$ &
		\textbf{B.}\quad $x=0,45$ &
		\textbf{C.}\quad $x=0,65$ &
		\textbf{D.}\quad $x=0,60$\\
	\end{tblr}

\subsection*{Question 11}
	Lequel de ces nombres est supérieur à 0,1 ?

	\medskip
	\begin{tblr}{colspec = {*4{X}}}
		\textbf{A.}\quad $\dfrac{1}{12}$ &
		\textbf{B.}\quad $10^{-2}$ &
		\textbf{C.}\quad $\dfrac{1}{5}$ &
		\textbf{D.}\quad $7\,\%$\\
	\end{tblr}

\subsection*{Question 12}
	On considère \quad $A=3 \times \dfrac{2}{5}+1$.\quad On a :

	\medskip
	\begin{tblr}{colspec = {*4{X}}}
		\textbf{A.}\quad $A=\dfrac{11}{5}$ &
		\textbf{B.}\quad $A=\dfrac{7}{15}$ &
		\textbf{C.}\quad $A=\dfrac{11}{10}$ &
		\textbf{D.}\quad $A=\dfrac{11}{15}$
	\end{tblr}


\needspace{10\baselineskip}
